\scanpage{149}
\mspage{94}
l'autre.

Du côté de $\mathcal M_{/S}$ : on doit comparer $X$ et $\delta\gamma(X)=C(X)$, c'est fait par
\[
i_X:X\longrightarrow C(X)
\]
qui est dans $\unsure{W_S|W}$, d'où l'isom. dans le localisé :
\[
X\simeq\delta\gamma(X)\quad\text{dans }(\mathcal M_{/S})W_S^{-1}\text{.}
\]
Il faut vérifier que c'est fonctoriel en $X$.

Du côté de $\unsure{\mathcal F|F}$ \deleted{$\underline{\mathrm{Fib}}\,S$}, comparer $X$ et $\gamma\delta(X)$, i.e. $X$ et $C(X)$ : Même tabac, on choisit $i_X:X\simeq C(X)=\gamma\delta(X)$. À vérifier la fonctorialité.
\figure{\begin{tikzpicture}[x=1cm,y=1cm]\draw (0,.16) .. controls (.02,.10) and (.09,.04) .. (.14,0);\end{tikzpicture}}

Il s'agit maintenant d'introduire un localisé plus fin de $\mathcal M_{/S}$. Pour ceci, on suppose \unsure{donné|donnée} une autre classe de morphismes, je l'appelle $\mathcal L$ (comme ``localisation'')
\[
\mathcal L\subset\mathrm{Fl}(\mathcal M)\text{,}
\]
p. ex. dans le cas de $\mathcal M=\mathrm{Cat}$, $\mathcal L$ \deleted{\unsure{sera}} sera formé des \deleted{\unsure{co}} isom. locaux$^{*}$ (s'il s'agit de définir $W_S^{\mathrm g}$), ou des \unsure{coisom.} locaux$^{**}$ (s'il s'agit de définir $W_S^{\mathrm d}$). On suppose, \inserted{que $\mathcal L$ contient les iso, et} $\mathcal L$ stable par composition, \deleted{\unsure{($^{***}$) qu'il contient les iso}}\,\unsure{,} mais non par chgt. de base.
% slanted
\margin{($^{*}$) ou aussi, des morphismes \uline{lisses}. (ça doit donner le $\eqsim$ résultat)}
% slanted
\margin{($^{**}$) ou aussi, des morphismes de projection}
% slanted
\margin{$^{***}$ Il n'y en \unsure{est|ait} peut-être pas besoin.}
On définit alors
\[
W_S^{\mathcal L}\subset\mathrm{Fl}(\mathcal M_{/S})
\]
par
\[
f\in W_S^{\mathcal L}\overset{\mathrm{def}}{\Longleftrightarrow}(f\times_S S')\in W,\quad\forall\,(S'\longrightarrow S)\in\mathcal L\text{.}
\]
% archive sheet number: 148
